\section{Related Work}
\label{sec:related}

\paragraph{Command-sequence CAD generation.}
The first wave of programmatic CAD generation operates on \emph{command-sequence representations} of sketch-and-extrude operations. DeepCAD~\citep{Wu2021} introduces a 178K-model corpus of Onshape parts encoded as discrete operation tokens, and Text2CAD~\citep{Khan2024} extends DeepCAD with 660K natural-language annotations spanning four abstraction levels, training a BERT-conditioned autoregressive transformer to emit those tokens. CAD-Coder~\citep{Guan2025} reformulates the same underlying data as CadQuery Python source and applies SFT + GRPO with a chain-of-thought prefix and Chamfer-distance reward, reaching mean CD $6.54\!\times\!10^{-3}$ on the Text2CAD test split.

\paragraph{The CadQuery-VLM lineage.}
A parallel research program has driven the recent state of the art in CadQuery code generation through three iterative releases. CAD-Recode~\citep{Rukhovich2025} first established the CadQuery formulation for reverse engineering, fine-tuning a Qwen2-1.5B language model on a procedurally generated corpus of 1M sketch-and-extrude scripts and reporting mean Chamfer distance 0.30 and IoU 92.0 on DeepCAD. cadrille~\citep{Kolodiazhnyi2026} extends this base by upgrading to Qwen2-VL-2B and unifying point-cloud, multi-view image, and text inputs in a single model, then introducing online reinforcement learning via Dr.CPPO with a piecewise IoU reward; the resulting cadrille-RL reaches DeepCAD image IoU 92.2 and Fusion360 image IoU 84.6 with 0.0\% invalidity. CADEvolve~\citep{Elistratov2026} shares the same backbone and RL recipe but expands the training corpus to 2.7M scripts via a GPT-5-mini-driven evolutionary loop, achieving DeepCAD image IoU 92.6. All three release model weights and inference scripts under permissive licences; we evaluate the latter two on BenchCAD as transfer baselines (\S\ref{sec:transfer}). These works collectively demonstrate the viability of CadQuery as a generation target and establish $\sim$92\% IoU as the saturated ceiling on existing benchmarks. They share three structural limitations BenchCAD addresses: training data is dominated by sketch+extrude operations, evaluation lacks both family-level taxonomy and standard-table grounding, and none decouple sub-capabilities --- every prior benchmark scores a single end-to-end metric, conflating visual perception, geometric abstraction, and code synthesis.

\paragraph{CAD code-generation and editing benchmarks.}
CADCodeVerify~\citep{Alrashedy2025} introduces \textbf{CADPrompt}, the first quantitative CAD code-generation benchmark --- 200 natural-language prompts paired with expert CadQuery code, scored via point-cloud distance and bounding-box IoU after ICP alignment. Despite establishing a useful protocol, CADPrompt is two orders of magnitude smaller than BenchCAD's verified split (20{,}143), lacks family taxonomy or difficulty stratification, leaks scale via absolute millimetre cues in prompts, and uses axis-aligned IoU that penalises geometrically correct but axis-permuted solutions. Beyond CADPrompt, two recent works target \emph{editing}: CAD-Editor proposes a locate-then-infill pipeline trained entirely on synthetic edit data without releasing a held-out human-verified edit benchmark, and CADialogue evaluates ten edit tasks within a conversational interface but does not release the eval set publicly. HistCAD~\citep{Dong2026} measures the preservation of ten geometric-constraint types (coincident, parallel, tangent) under edits via a FreeCAD edit-then-check protocol, but operates at the sketch-constraint level rather than at the executable-program level. To our knowledge, BenchCAD-Edit is the first public CAD benchmark with (i) execution-verified before/after CadQuery code pairs, (ii) a non-target preservation metric distinguishing target-feature success from program-wide locality, and (iii) cross-family balanced coverage with explicit edit-category taxonomy.

\paragraph{Vision-language spatial and code benchmarks.}
Recent benchmark methodology in adjacent domains informs our design. MMSI-Bench~\citep{MMSIBench2026} evaluates 37 multimodal LLMs on 1{,}000 hand-curated multi-image spatial reasoning questions across an 11-task taxonomy, reporting a 67-point human--model gap and showing that scaling model size yields negligible gains. AutoCodeBench~\citep{Chou2026} constructs a 3{,}920-problem multilingual code-generation benchmark via an automated LLM--sandbox pipeline and releases Lite/Complete subsets for differentiated evaluation regimes. Infinity-Chat~\citep{Hivemind2025} curates 26K real-world open-ended queries with 31K dense human annotations and reveals an ``Artificial Hivemind'' homogenisation effect across 70+ models. We adopt three of their structural patterns: explicit task taxonomy with per-task representatives (MMSI-Bench), Lite + Full split design for differentiated evaluation cost (AutoCodeBench), and a named, quantitative finding as the paper's punch line.

\paragraph{Position of BenchCAD.}
BenchCAD occupies a vacant intersection: a CAD code-generation benchmark that is simultaneously (i) \emph{execution-verified at scale} (20{,}143 parts at IoU $\geq 0.99$), (ii) \emph{standard-anchored} (55\% of families bound to ISO/DIN/EN/ASME/IEC specification tables), (iii) \emph{operation-rich} (45+ CadQuery operations versus 2--5 in prior corpora), and (iv) \emph{capability-decomposed} (five tasks isolating visual reconstruction, parametric abstraction, and code synthesis along the image$\to$code causal chain, including the first verified parametric edit task). Table~\ref{tab:cmp_prior} summarises this comparison.

\begin{table*}[t]
\centering
\small
\setlength{\tabcolsep}{5pt}
\caption{\textbf{BenchCAD versus prior CAD code-generation work.} BenchCAD is a \emph{unified, capability-decomposed evaluation framework} for large-model CAD reasoning, distinct in purpose from the CadQuery-VLM lineage (CAD-Recode, cadrille, CADEvolve) which releases training corpora alongside small fine-tuned models (Qwen 1.5B--2B) saturating $\sim$92\% IoU on legacy sketch+extrude splits. \textbf{Type:} \emph{C+M} = corpus+model, \emph{D} = dataset, \emph{B} = benchmark. \textbf{\#Fam.} = named industrial part families with explicit functional identity. \textbf{\#Std} = distinct ISO/DIN/EN/ASME/IEC specification codes. \textbf{\#Ops} = distinct CadQuery API methods exercised in the released corpus (static analysis of every released \texttt{.py}; see Appendix~\ref{app:ops}). \textbf{Adv} = all four advanced operation families (helix, loft/sweep, twist-extrude, parametric involute-gear) exercised. $\checkmark$ = supported, p.\ = partial, -- = absent.}
\label{tab:cmp_prior}
\begin{tabular}{lccrrrrccc}
\toprule
Benchmark & Yr & Type & \#Sam. & \#Fam. & \#Std & \#Ops & Adv & Edit & QA \\
\midrule
DeepCAD~\citep{Wu2021}                   & '21 & C+M & 178K  & -- & 0  & 4   & --   & --   & --   \\
Fusion360 G.~\citep{Willis2021fusion}    & '21 & D   & 8K    & -- & 0  & 6   & --   & --   & --   \\
Text2CAD~\citep{Khan2024}                & '24 & C+M & 170K  & -- & 0  & 17  & --   & --   & --   \\
CADPrompt~\citep{Alrashedy2025}          & '25 & B   & 200   & -- & 0  & --  & p.\  & --   & --   \\
CAD-Coder~\citep{Guan2025}               & '25 & C+M & 110K  & -- & 0  & 4   & --   & p.\  & --   \\
CAD-Recode~\citep{Rukhovich2025}         & '25 & C+M & 1M+7K & -- & 0  & 18  & --   & p.\  & seq.\\
cadrille~\citep{Kolodiazhnyi2026}        & '26 & C+M & 1M+7K & -- & 0  & 18  & --   & --   & --   \\
CADEvolve~\citep{Elistratov2026}         & '26 & C+M & 2.7M  & -- & 0  & 25  & p.\  & --   & --   \\
\midrule
\textbf{BenchCAD (ours)} & \textbf{'26} & \textbf{D+B} & \textbf{17{,}900} & \textbf{106} & \textbf{43} & \textbf{49} & $\checkmark$ & $\checkmark$ & $\checkmark$ \\
\bottomrule
\end{tabular}
\end{table*}
